\documentclass[sigconf]{acmart}
\AtBeginDocument{%
  \providecommand\BibTeX{{%
    Bib\TeX}}}
    
\usepackage{kotex} %한글사용
\usepackage{lipsum}
\usepackage{algorithm}
\usepackage{algpseudocode}

\setcopyright{acmlicensed}
\copyrightyear{2018}
\acmYear{2018}
\acmDOI{XXXXXXX.XXXXXXX}
\acmConference[Conference acronym 'XX]{Make sure to enter the correct
  conference title from your rights confirmation email}{June 03--05,
  2018}{Woodstock, NY}

\acmISBN{978-1-4503-XXXX-X/2018/06}

\begin{document}

\title{CSE331 - Assignment \#1}

\author{\textcolor{blue}{Hojun Lee (20211277)}}
\affiliation{%
  \institution{UNIST}
  \country{South Korea}
}
\email{hojunlee@unist.ac.kr}

\renewcommand{\shortauthors}{\textcolor{blue}{Kim et al.}}

\settopmatter{printacmref=false} % Removes citation block in the first page
\renewcommand\footnotetextcopyrightpermission[1]{} % Removes footnote with conference info

\maketitle

\section{PROBLEM STATEMENT}
\section{BASIC SORTING ALGORITHMS}
\subsection{Bubble Sort}
\subsubsection{Original Bubble Sort}

버블 정렬(Bubble Sort)은 가장 직관적인 정렬 알고리즘 중 하나입니다. 이 알고리즘은 배열을 반복적으로 순회하며, 현재 원소와 다음 원소를 비교합니다. 만약 현재 원소가 다음 원소보다 크다면, 두 원소의 위치를 교환합니다. 만약 다음 원소가 더 작다면, 두 원소의 위치를 바꾼다. 이 알고리즘은 배열을 1번 탐색할 떄, 가장 큰 원소를 오른쪽 끝에 위치하게한다. 떄문에 총 비교 회수는 n-1 + n-2 + ... + 1 = n(n-1)/2 번으로 시간복잡도 는 n2 이다.

\subsubsection{optimized bubble sort}
모든 원소가 정렬 되어있다면, 더이상 비교할 필요가 없다.
배열 이 정렬되기 위해, 이번 탐색에 교환이 일어났는지를 확인하는
is\_swap 

불리언 변수를 이용한다.
교환이 일어나지 않았다면 배열이 정렬됨을 의미한다.
is\_swap 은 false 값을 가지고, 정렬이 종료된다.
이 보고서에서 bubble sort는 이 최적화를 이용한 알고리즘을 의미한다.
이 경우 시간 복잡도는 최악 평균 은 n2으로 최적화 하지 않은 경우와 같지만.
모든 배열이 정렬 된 경우 최선의 시간복잡도를 가진다.
이 경우 1번의 순회에 정렬이 종료되어, 시간복잡도는 O(n) 이다.

이 알고리즘은 구현이 단순하고 직관적이다. 또한 공간복잡도는 O(n) 으로 in-place 하고 stable 한 정렬 알고리즘이다.

하지만, 시간복잡도가 최악 n2, 기대값이 n2 으로 정렬해야하는 원소 수 가 많아지면, 성능이 떨어질 수 있다. 또한, 불필요한 교환이 일어난다.
\subsection{Selection Sort}
알고리즘 설명
장점: 메모리는 적게쓴다.
단점: 항상 n2 만큼 걸려.

\subsection{Insertion Sort}
online algorithm
in-place algorithm
새로운 원소를 받았을 떄, O(n) 에 정렬가능

\subsubsection{binary search}
\subsection{Merge Sort}
분할 정복 패러다임을 이용한다. theta(nlogn) 만큼 걸림
재귀적인 동작방식
not inplace -> n 개 정렬하는데 n개의 원소를 저장하는 추가적인 공간이 필요함.
locality 괜춘혀만족
stable 함
재귀와 locality 땜에 오버헤드가 크다. 떄문에 어떤 작은 입력크기 n 이하에서는 위의 정렬이 더 빠른 구간이 생김.
공간복잡도 n




\subsection{Heap Sort}
heap 을 구성하고 heap 에서 1원소를 뺀다. 이를 반복해, 정렬한다. heap을 구성하는데 O(n) 의 시간이 걸린다. 그리고 heap 에서 원소를 뺴고, heap 구조를 복원하는데 O(logn) 만큼 걸린다. 이를 n번 반복하면 heapsort의 시간 복잡도는 O(nlogn) 이다. 어떤경우에도 같은 방식으로 동작하기 떄문에 theta(nlogn) 이다.
heap에서 추출된 값들을 저장해둘 임시 값이 필요하므로 추가 공간은 O(n) 이다. 따라서 in-place 한 알고리즘은 아니다.
어떠한 경우에도 O(nlogn) 이란 시간이 걸리는 점은 아쉬운 부분이다. 또한, array base heap 에서 노드와 자식노드와 비교할떄, 연속적인 메모리 공간에 위치하지 않으므로 locality 가 떨어진다.

\subsection{quick Sort}
\subsubsection{양 끝 중 하나를 피벗으로 하는 quick sort}
분할 정복 패러다임을 이용한 정렬 알고리즘이다. 피벗을 잡고 피벗의 왼쪽에는 피벗보다 작거나 같은 원소가, 오른쪽에는 피벗 보다 큰원소가 위치하게 한다. 피벗의 왼쪽 배열과 피벗의 오른쪽 배열에 대해 각각 재귀적으로 반복한다. best case 는  양 배열을 균등하게 나누는 경우이고 이 경우 O(nlogn) 만큼의 시간 복잡도를 가진다. worst case 의 경우에는 피벗을 제외한 나머지가 한쪽에 몰려있는 경우이고, 이경우에 O(n2) 의 시간복잡도를 가지게 된다.
quick sort는 원소를 저장하기위한 큰 추가적인 메모리를 필요로 하지 않는다. 하지만, 재귀적 동작을 한다. 그러므로 best case 에서 양쪽을 반으로 나누면 추가로 필요한 메모리는 O(1)이고 O(logn) 번의 깊이를 가지므로 O(logn) 이 추가로 필요하다. 따라서 in-place 알고리즘이다.
피벗을 어떻게 정해지는지에 따라서 성능이 좌우되는 알고리즘이다. 피벗을 처음이나 가장 오른쪽을 정하는 경우, 정렬된 배열에서 worst case 와 비슷해지므로, 성능이 안좋다. 떄문에 범위내 random 하게 피벗을 정하는 알고리즘을 사용하였다.

\section{ADVANCED SORTING ALGORITHMS}
\subsection{Library Sort}
\subsection{Comb Sort}
\subsection{Cocktail Shaker Sort}
\subsection{Introsort}
\subsection{Tournament Sort}
\subsection{Tim Sort}
\section{EXPERIMENTAL RESULTS AND ANALYSIS}

% \begin{algorithm}[H]
% \caption{Bubble Sort Algorithm}
% \begin{algorithmic}[1] % [1] 옵션을 추가하여 줄 번호 활성화
% \Procedure{BubbleSort}{A}
%     \State $n \gets \text{length of } A$
%     \Repeat
%         \State $swapped \gets \textbf{false}$
%         \For{$i \gets 1$ to $n-1$}
%             \If{$A[i-1] > A[i]$}
%                 \State Swap $A[i-1]$ and $A[i]$
%                 \State $swapped \gets \textbf{true}$
%             \EndIf
%         \EndFor
%         \State $n \gets n - 1$
%     \Until{$\textbf{not } swapped$}
% \EndProcedure
% \end{algorithmic}
% \end{algorithm}

\bibliographystyle{ACM-Reference-Format}
\bibliography{reference}

\end{document}
